\documentclass[12pt]{article}
\usepackage{amsmath, amssymb, geometry}
\geometry{margin=1in}
\setlength{\parskip}{0.5em}
\setlength{\parindent}{0pt}

\begin{document}

\title{\textbf{Solution Key}\\
Labour Economics MT2}
\date{}
\maketitle

\section*{Multiple Choice Solutions with Explanations}

\begin{enumerate}

\item \textbf{B} \\
A monopolist faces marginal revenue below price, so labour demand is based on marginal revenue product (MRP), which is less elastic.



\textbf{2.A}

\subsection*{Economic Intuition}

Economic rent is defined as the difference between the wage a worker receives and the worker’s \textbf{reservation wage} (the minimum wage required to induce them to work).

\[
\text{Economic Rent} = W - W^{R}
\]

where:
\begin{itemize}
    \item \(W\) = actual wage
    \item \(W^{R}\) = reservation wage
\end{itemize}


If workers are earning economic rents, then:

\begin{itemize}
    \item They are being paid \textbf{more than necessary} to keep them employed.
    \item A reduction in wages \textbf{will not immediately reduce employment}.
    \item Workers will continue working as long as:
    \[
    W \geq W^{R}
    \]
\end{itemize}

\subsection*{Why Answer A is Correct}

Answer A reflects the idea that:

\begin{quote}
A firm can reduce wages \textbf{without losing workers}, as long as wages remain above the reservation wage.
\end{quote}

This is the defining property of economic rents:
\begin{itemize}
    \item Firms have some flexibility to lower wages.
    \item Workers do not leave immediately because they are still better off than their next best alternative.
\end{itemize}

\subsection*{Why Answer D is Incorrect}

Answer D states that workers would leave if wages are reduced. This is incorrect because:

\begin{itemize}
    \item Workers only leave when:
    \[
    W < W^{R}
    \]
    \item A small wage reduction \textbf{within the rent margin} does not trigger exit.
\end{itemize}

Thus, Answer D ignores the key concept of a \textbf{reservation wage threshold}.

\subsection*{Conclusion}

\[
\boxed{
\text{Economic rents imply that wages can fall without reducing employment, as long as } W \geq W^{R}.
}
\]

Therefore, \textbf{Answer A is correct}.

\item \textbf{B} \\
In monopsony, marginal cost of labour exceeds the wage due to upward-sloping supply.

\item \textbf{A} \\
Industry labour supply is generally upward sloping because higher wages attract more workers.

\item \textbf{A} \\
A payroll tax raises labour costs and reduces employment.

\item \textbf{C} \\
Tax incidence depends on relative elasticities, so the burden is shared.

\item \textbf{D} \\
Elasticity is percentage change in labour demand over percentage change in wages.

\item \textbf{E} \\
Isoprofit curves are concave due to diminishing marginal returns to safety investment.

\item \textbf{D} \\
Efficiency wages operate through morale, effort, and turnover—not directly through safety.

\item \textbf{B} \\
Indifference curves reflect diminishing marginal rate of substitution.

\item \textbf{A} \\
More risk-averse workers require higher compensating differentials.

\item \textbf{A} \\
Higher isoprofit curves represent lower profits.

\item \textbf{B} \\
The wage-safety locus shows the trade-off between wages and job risk.

\item \textbf{E} \\
Adam Smith introduced compensating wage differentials.

\item \textbf{B} \\
Doctors respond to incentives; wage differentials can induce relocation.

\item \textbf{D} \\
Competitive equilibrium assumes perfect information and mobility.

\item \textbf{A} \\
Risk-averse workers match with safer firms.

\item \textbf{E} \\
Uber drivers accept lower wages for flexibility.

\item \textbf{B} \\
Differentials compensate for job characteristics, not utility directly.

\item \textbf{D} \\
Signalling theory focuses on screening rather than productivity.

\item \textbf{B} \\
Human capital increases productivity and earnings.

\item \textbf{C} \\
Opportunity cost is forgone earnings while studying.

\item \textbf{C} \\
Private costs/benefits accrue to the individual.

\item \textbf{E} \\
All statements describe age-earnings profiles correctly.

\item \textbf{E} \\
With perfect capital markets, individuals do not rely solely on current income.

\item \textbf{A} \\
Optimal schooling equates marginal benefit and marginal cost.

\item \textbf{A} \\
Imperfect information exists when key variables are unknown.

\item \textbf{C} \\
Earnings rise until mid-career, then decline.

\item \textbf{A} \\
$PV \approx \frac{Y}{r}$ applies to long-term income streams.

\item \textbf{B} \\
Monopsony equilibrium: $MCL = VMP \Rightarrow (W_M, N_M)$.

\item \textbf{C} \\
Competitive equilibrium: labour supply equals labour demand.

\item \textbf{A} \\
Monopsonist sets optimal employment directly, so no vacancies.

\item \textbf{A} \\
Area $d$ represents direct schooling costs.

\item \textbf{D} \\
Area $e$ is the earnings premium from higher education.

\item \textbf{E} \\
Areas $b + c + d$ represent total education costs.

\end{enumerate}

\section*{Final Answer Summary}

\[
\begin{aligned}
&1.\,B \quad 2.\,D \quad 3.\,B \quad 4.\,A \quad 5.\,A \\
&6.\,C \quad 7.\,D \quad 8.\,E \quad 9.\,D \quad 10.\,B \\
&11.\,A \quad 12.\,A \quad 13.\,B \quad 14.\,E \quad 15.\,B \\
&16.\,D \quad 17.\,A \quad 18.\,E \quad 19.\,B \quad 20.\,D \\
&21.\,B \quad 22.\,C \quad 23.\,C \quad 24.\,E \quad 25.\,E \\
&26.\,A \quad 27.\,A \quad 28.\,C \quad 29.\,A \\
&30.\,B \quad 31.\,C \quad 32.\,A \quad 33.\,A \quad 34.\,D \quad 35.\,E
\end{aligned}
\]

\end{document}